\section{灵敏度分析}

\textbf{灵敏度分析}是通过改变模型中的关键参数或输入数据，观察模型输出及最终结论的变化程度，从而判断模型对参数扰动是否敏感、结果是否稳定；\textbf{误差分析}是比较模型计算结果与真实值、观测值或参考值之间的差异，并分析误差大小及其来源，从而评价模型的准确性。二者的主要区别在于：灵敏度分析关注“\textbf{参数变化会使结果改变多少}”，主要检验模型的稳定性和鲁棒性；误差分析关注“\textbf{模型结果与实际值相差多少}”，主要检验模型的准确性和可信度。

\subsection{误差分析}
\begin{enumerate}[label=(\arabic*), leftmargin=2.5em, itemsep=0.3em]
\item \textbf{检验模型可靠性。} 模型求解后，应分析参数变化、数据误差和模型假设对结果的影响，判断结论是否稳定可信。

\item \textbf{选择合适的检验方法。} 并非所有模型都必须同时进行多种检验，应根据模型类型和数据特点选择灵敏度分析、误差分析或其他检验方法。

\item \textbf{围绕关键结果展开。} 重点检验重要参数、核心指标和最终方案，不必对所有变量逐一分析。

\item \textbf{结果应有解释。} 不能只给出检验图表或误差数值，还应说明模型是否稳定、误差是否可接受以及主要影响因素是什么。
\end{enumerate}
\hspace*{\parindent} \textbf{模型检验的目的，是说明结果是否稳定、准确、可信，而不是简单增加图表和计算。}
%\hspace*{\parindent} 默认系统缩进

\textbf{一、灵敏度分析}
\begin{enumerate}[label=(\arabic*), leftmargin=2.5em, itemsep=0.3em]
\item \textbf{选择关键参数。} 优先分析取值具有不确定性、对模型结构影响较大或来源于人为设定的参数。

\item \textbf{合理设置变化范围。} 可在基准值附近按一定比例调整参数，如分别变化 (±5\%)、(±10\%) 或其他符合实际的范围。

\item \textbf{保持其他条件不变。} 每次改变一个参数，观察目标值、评价结果、排序或最优方案的变化；必要时可分析多个参数的共同影响。

\item \textbf{展示变化规律。} 可使用表格、折线图或柱状图表示参数变化与模型结果之间的关系。

\item \textbf{判断模型稳定性。} 若参数小幅变化时结果变化较小，说明模型稳定；若结果变化明显，应指出敏感参数并分析原因。
\end{enumerate}
\hspace*{\parindent} \textbf{灵敏度分析重点回答：参数发生变化时，模型结论是否仍然成立。}

\textbf{二、误差分析}
\begin{enumerate}[label=(\arabic*), leftmargin=2.5em, itemsep=0.3em]
\item \textbf{明确误差来源。} 误差可能来自原始数据、测量过程、缺失值处理、参数估计、模型简化或数值计算。

\item \textbf{选择误差指标。} 可根据问题类型采用绝对误差、相对误差、均方误差、平均绝对误差、拟合优度等指标。

\item \textbf{进行结果对比。} 将模型结果与真实值、历史数据、题目给定结果或其他合理方法的结果进行比较。

\item \textbf{分析误差大小。} 说明误差是否处于可接受范围，并结合实际背景解释偏差产生的原因。

\item \textbf{指出模型局限。} 对误差较大的情况，应说明模型在哪些条件下可能失效，以及可以如何改进。
\end{enumerate}
\hspace*{\parindent} \textbf{误差分析重点回答：模型结果与实际情况相差多少，误差从哪里产生。}

\textbf{写作提示}
\begin{enumerate}[label=(\arabic*), leftmargin=2.5em, itemsep=0.3em]
\item 灵敏度分析和误差分析应与具体模型结合，避免只介绍方法而没有实际计算。

\item 图表应突出关键变化和主要结论，不必展示大量重复结果。

\item 不能仅凭一次计算就说明模型稳定，应设置合理的参数变化范围或对照数据。

\item 若缺少真实数据，无法直接计算误差，应如实说明，并采用交叉验证、历史回代或情景分析等方式间接检验。
\end{enumerate}
\hspace*{\parindent} \textbf{写作原则：指标合理、过程清楚、结论明确、评价客观。}
